% =====================================================================
%  CUERPO DEL MANUAL DE DESARROLLO - PharmaWeigh
%  Capturas reales del sistema. Sin credenciales.
% =====================================================================

\seccion{1. Acceso al sistema}

\tablaAcceso

\vspace{8pt}

\reglasAcceso

\captura{01_login.png}{Pantalla de entrada.}

\avisoValidacion

% =====================================================================
\newpage
\seccion{2. Tu papel dentro del sistema}

Como \manualRol{} cargas las fórmulas de producción. Lo que guardes en una receta es,
literalmente, lo que va a aparecer en la pantalla del operario y lo que el sistema va
a exigir en la balanza.

\begin{center}
\begin{tabularx}{\linewidth}{@{}>{\raggedright\arraybackslash}X c@{}}
\toprule
\textbf{Acción} & \textbf{Puedes} \\
\midrule
\textbf{Crear recetas}                              & Sí \\
Consultar órdenes y el detalle de cada pesaje       & Sí \\
Consultar inventario, materiales y lotes            & Sí \\
Consultar la bitácora de auditoría                  & Sí \\
\midrule
\textbf{Editar o borrar una receta}                 & \textbf{No (nadie puede)} \\
Crear órdenes de producción                         & \textbf{No} \\
Dar de alta materiales                              & \textbf{No} \\
Recibir lotes o cambiar su estado                   & \textbf{No} \\
Registrar pesajes                                   & \textbf{No} \\
Firmar fases                                        & \textbf{No} \\
Ver usuarios o códigos de barras                    & \textbf{No} \\
\bottomrule
\end{tabularx}
\end{center}

\subseccion{Tu menú}

\menu{Dashboard}, \menu{Órdenes}, \menu{Dispensado}, \menu{Recetas},
\menu{Inventario} y \menu{Auditoría}.

\captura{03_dashboard.png}{Dashboard: conteos del día, procesos activos y actividad reciente.}

% =====================================================================
\newpage
\seccion{3. Cómo está armada una receta}

\vspace{6pt}
\begin{center}
\begin{tikzpicture}[
    every node/.style={font=\footnotesize},
    caja/.style={draw=pharmablue, fill=pharmablue!6, rounded corners=5pt,
        minimum width=3.2cm, minimum height=0.8cm, align=center},
    sub/.style={draw=pharmagreen, fill=pharmagreen!6, rounded corners=5pt,
        minimum width=3.2cm, minimum height=0.7cm, align=center},
    l/.style={-{Stealth[length=5pt]}, thick, color=gristexto}
]
\node[caja] (r) {\textbf{Receta}\\ \tiny código, nombre, rendimiento};
\node[caja, below=0.5cm of r] (f1) {\textbf{Fase 1} \tiny(nombre + instrucciones)};
\node[sub, below=0.4cm of f1] (i1) {Ingrediente 1 \tiny(material, cantidad, tolerancias)};
\node[sub, below=0.25cm of i1] (i2) {Ingrediente 2};
\node[caja, below=0.5cm of i2] (f2) {\textbf{Fase 2}};
\node[sub, below=0.4cm of f2] (i3) {Ingrediente 3};
\draw[l] (r) -- (f1);
\draw[l] (f1) -- (i1);
\draw[l] (i2) -- (f2);
\draw[l] (f2) -- (i3);
\end{tikzpicture}
\end{center}
\vspace{6pt}

Tres niveles: la receta contiene \textbf{fases}, y cada fase contiene
\textbf{ingredientes}. Eso no es decorativo:

\begin{itemize}
    \item El \textbf{orden de las fases} y el \textbf{orden de los ingredientes dentro
          de cada fase} es el orden exacto en que el operario tiene que pesar.
          \textbf{No lo puede alterar}: el servidor lo impone.
    \item \textbf{Cada fase termina en una firma electrónica} de Supervisor, Calidad o
          Administrador, y \textbf{bloquea la fase siguiente} hasta que se firme.
\end{itemize}

\begin{cajaOjo}
\textbf{Cada fase que agregas es una firma que alguien tiene que ir a dar a la
estación.} Divide por etapas reales del proceso (dispensado de activos, mezclado,
granulación), no por comodidad de captura. Tres fases son tres interrupciones de la
línea esperando a un supervisor.
\end{cajaOjo}

% =====================================================================
\newpage
\seccion{4. Crear una receta}

\subseccion{4.1 Paso a paso}

\begin{enumerate}
    \item \menu{Recetas} $\rightarrow$ \boton{+ Nueva Receta}. El formulario se abre en
          la misma pantalla.
    \item Datos generales:
        \begin{itemize}
            \item \campo{Código}: identificador único (ej. \codigo{REC-PCT-500}).
            \item \campo{Nombre}: descriptivo (ej. \textit{Tableta Paracetamol 500mg}).
            \item \campo{Rendimiento} y \campo{Unidad Rendimiento}: kg, g, L, unidades
                  o tabletas.
            \item \campo{Descripción}: opcional.
        \end{itemize}
    \item \boton{+ Agregar Fase} por cada etapa. Cada fase pide su nombre (obligatorio)
          y unas \campo{Instrucciones de la fase} opcionales, que el operario verá en la
          cabecera mientras trabaja esa fase.
    \item Dentro de cada fase, \boton{+ Agregar Material} por cada ingrediente:
        \begin{itemize}
            \item \campo{Material}: de la lista de materiales ya dados de alta.
                  Entre paréntesis viene su unidad.
            \item \campo{Cantidad}: el target para \textbf{un} lote estándar.
            \item \campo{Tol -} y \campo{Tol +}: porcentajes (por omisión \cd{-2} y \cd{2}).
            \item \campo{Instrucciones}: el texto que el operario lee en ese paso.
            \item \campo{Peligroso}: enciende la alerta roja en la estación.
        \end{itemize}
    \item \boton{Guardar Receta}.
\end{enumerate}

\captura{13_receta_nueva.png}{Formulario de receta: primero las fases, y dentro de cada una sus materiales.}

\vspace{6pt}

El botón \boton{Guardar Receta} permanece \textbf{gris y deshabilitado} mientras no
haya al menos una fase y todas las fases tengan al menos un material. Si el servidor
rechaza el alta, el mensaje rojo aparece arriba del formulario y \textbf{no se guarda nada}.

\begin{cajaOjo}
\textbf{Una receta guardada no se edita, no se versiona y no se borra.} No existe esa
pantalla, ni para ti ni para nadie. Revisa códigos, cantidades, unidades y tolerancias
\textbf{antes} de pulsar el botón. Si te equivocaste, la única salida es dar de alta
otra receta con código distinto y coordinar que las órdenes nuevas usen esa.
\end{cajaOjo}

\newpage
\subseccion{4.2 Materiales: el prerrequisito}

Sólo puedes elegir materiales que ya existan. \textbf{Tu rol no los da de alta}: eso lo
hacen Almacén, Supervisión o el Administrador desde \menu{Inventario}.

\captura{14_inventario_materiales.png}{Materiales disponibles, con su código y su unidad.}

\vspace{6pt}

Antes de empezar una receta, revisa esta tabla y confirma dos cosas:

\begin{enumerate}
    \item Que existan todos los materiales que vas a necesitar. Si falta alguno,
          pídeselo a Almacén antes de abrir el formulario: si lo das de alta a media
          captura, pierdes lo que llevabas escrito.
    \item Que la \textbf{unidad} de cada material sea la correcta. La cantidad que
          teclees se interpreta en \textit{esa} unidad, y \textbf{la unidad de un
          material no se puede cambiar después}. Un material dado de alta en gramos
          cuando debía ser kilogramos convierte todos tus targets en un error de
          factor 1000.
\end{enumerate}

% =====================================================================
\newpage
\seccion{5. Tolerancias: cómo las calcula el sistema}

\subseccion{5.1 La fórmula}

Las tolerancias son \textbf{porcentaje del target}, y el target real se obtiene
multiplicando la cantidad de la receta por el multiplicador de la orden.

\begin{center}
\begin{tabularx}{0.95\linewidth}{@{}l >{\raggedright\arraybackslash}X@{}}
\toprule
Target real & cantidad de la receta $\times$ multiplicador de la orden \\
Mínimo      & target $\times$ (1 + Tol$-$/100) \\
Máximo      & target $\times$ (1 + Tol$+$/100) \\
\bottomrule
\end{tabularx}
\end{center}

\vspace{6pt}

Ejemplo real del sistema: 5.000 kg con \cd{-1} y \cd{1} da un rango de
\textbf{4.950 a 5.050 kg}. En una orden con multiplicador 2, el mismo ingrediente pide
10.000 kg con rango 9.900 a 10.100 kg: \textbf{la tolerancia escala junto con el target}.

\begin{cajaTip}[Precisión]
El sistema trabaja con \textbf{4 decimales} y los porcentajes admiten centésimas
(\cd{-2.5} es válido). Todo redondeo va \textbf{hacia adentro} del rango: un decimal
sobrante nunca ensancha la tolerancia. Un target muy pequeño con una tolerancia muy
cerrada puede quedar en un rango de un par de milésimas, imposible de acertar en la
balanza de piso. Verifica que el rango que sale sea pesable con el instrumento real.
\end{cajaTip}

\subseccion{5.2 El semáforo que verá el operario}

\begin{center}
\begin{tabularx}{0.9\linewidth}{@{}c l >{\raggedright\arraybackslash}X@{}}
\toprule
& \textbf{Estado} & \textbf{Cuándo} \\
\midrule
\semaverde{}    & \cd{OK}         & Dentro del rango y lejos de las orillas. \\
\semaamarillo{} & \cd{PRECAUCIÓN} & Dentro del rango, pero en el \textbf{10\% exterior} de la banda. Se puede confirmar. \\
\semarojo{}     & \cd{FUERA}      & Fuera del rango. \textbf{Botón bloqueado.} \\
\bottomrule
\end{tabularx}
\end{center}

\vspace{6pt}

La franja ámbar no la defines tú: el sistema la calcula sola como el 10\% exterior del
rango que pusiste. Con \cd{-1}/\cd{+1} sobre 5 kg, el ámbar cae en los 10 g pegados a
cada extremo.

\begin{cajaOjo}
\textbf{Un peso fuera de tolerancia no se registra}, y \textbf{no hay forma de
autorizar una excepción} desde ninguna pantalla. Si pones una tolerancia demasiado
cerrada, detienes la línea y la única salida es crear otra receta. Si la pones
demasiado ancha, el sistema deja pasar pesos que no debería: no tiene criterio propio,
respeta el tuyo.
\end{cajaOjo}

\subseccion{5.3 Criterios habituales}

$\pm 1\%$ para principios activos, $\pm 2\%$ a $\pm 3\%$ para excipientes, $\pm 5\%$
para lubricantes y deslizantes. Ajústalo al expediente de desarrollo de la planta, no
a este manual.

% =====================================================================
\newpage
\seccion{6. Instrucciones y materiales peligrosos}

\subseccion{6.1 Instrucciones}

Hay dos niveles, y los dos llegan al operario:

\begin{itemize}
    \item \campo{Instrucciones de la fase}: se muestran en la cabecera azul mientras dure
          esa fase completa. Sirven para condiciones generales (\textit{trabajar bajo
          campana}, \textit{temperatura máxima 25 °C}).
    \item \campo{Instrucciones} del ingrediente: se muestran en el recuadro del paso, junto
          al material. Sirven para la operación concreta (\textit{tamizar antes de pesar
          (malla 40)}, \textit{agregar al final}).
\end{itemize}

\begin{cajaTip}[Cómo escribirlas]
Frases cortas, en imperativo, con el dato numérico dentro. El operario las lee de pie,
con guantes y con prisa. \textit{``Tamizar malla 40 antes de pesar''} funciona;
\textit{``se recomienda considerar el tamizado previo''} no.
\end{cajaTip}

\subseccion{6.2 La casilla \campo{Peligroso}}

Marca la casilla y la estación pinta el recuadro del paso en rojo con la etiqueta
\textbf{MATERIAL PELIGROSO}. Es una advertencia visual: \textbf{no cambia ninguna regla
del sistema}, no exige confirmación extra ni restringe quién puede pesarlo. El equipo
de protección concreto ponlo en las instrucciones del ingrediente.

% =====================================================================
\newpage
\seccion{7. Verificar cómo quedó tu receta}

No hay una vista de detalle de receta en \menu{Recetas}: la tabla sólo muestra código,
nombre, versión, rendimiento, número de fases y de ingredientes.

\captura{12_recetas.png}{Listado de recetas.}

\vspace{6pt}

Para revisar targets y rangos tal como los va a ver la planta, pide a Supervisión que
cree una orden de prueba con tu receta y ábrela desde \menu{Órdenes}. Tu rol aterriza
en la \textbf{vista de consulta}, que muestra cada fase y, por ingrediente, el target y
el rango calculado.

\captura{24_orden_consulta.png}{Vista de consulta de una orden: por ingrediente, target, rango permitido y peso real dispensado.}

\vspace{6pt}

Es también la mejor forma de cerrar el ciclo: una vez producida la orden, la misma
pantalla te muestra el \textbf{peso real} que se dispensó de cada ingrediente. Si los
pesos reales se pegan sistemáticamente a una orilla del rango, la tolerancia o el
target merecen revisión en la siguiente fórmula.

\captura{04_ordenes.png}{Listado de órdenes. Para tu rol la acción siempre es \boton{Ver detalle}.}

% =====================================================================
\newpage
\seccion{8. Auditoría}

\menu{Auditoría} te da los 500 eventos más recientes, con filtros por acción y por
usuario. No hay filtro por fecha, ni búsqueda de texto, ni exportación.

\captura{17_auditoria.png}{Bitácora del sistema.}

\vspace{6pt}

Para tu trabajo, dos filtros valen la pena:

\begin{itemize}
    \item \cd{CREAR\_RECETA}: qué fórmulas se cargaron, por quién y cuándo.
    \item \cd{DISPENSAR}: los pesos reales registrados, útil para contrastar contra los
          targets que diseñaste.
\end{itemize}

La bitácora no se puede editar ni borrar desde la aplicación.

% =====================================================================
\newpage
\seccion{9. Qué sí puedes romper}

\begin{center}
\begin{tabularx}{\linewidth}{@{}p{0.32\linewidth} >{\raggedright\arraybackslash}X@{}}
\toprule
\textbf{Riesgo} & \textbf{Por qué importa} \\
\midrule
\textbf{Guardar una receta con un error} & No se edita ni se borra. El error se va a producción tal cual y sólo se resuelve creando otra receta. \\
\textbf{Confundir la unidad del material} & Si el material está en gramos y tú piensas en kilogramos, el target sale con un factor 1000 y nadie lo va a poder pesar. \\
\textbf{Tolerancia demasiado cerrada}    & Detienes la línea: el operario no logra confirmar y no existe autorización de excepción. \\
\textbf{Tolerancia demasiado ancha}      & El semáforo deja pasar pesos fuera de especificación. El sistema respeta lo que pusiste. \\
\textbf{Orden de ingredientes equivocado} & Es el orden real de pesaje y no se puede cambiar en piso. \\
\textbf{Demasiadas fases}                & Cada fase exige una firma presencial. Fragmentar de más paraliza la producción. \\
\textbf{Olvidar marcar \campo{Peligroso}} & El operario pierde la advertencia visual del material. \\
\bottomrule
\end{tabularx}
\end{center}

% =====================================================================
\newpage
\seccion{10. Preguntas frecuentes}

\subseccion{¿Puedo modificar una receta que ya tiene órdenes?}
No, y tampoco una que no las tenga. Ninguna receta se edita ni se borra en el sistema.
Crea otra con código distinto.

\subseccion{¿Cómo hago una versión 2 de una fórmula?}
No hay función de versionado. El campo \textit{versión} existe, pero ninguna pantalla
lo incrementa. En la práctica: nueva receta, código nuevo (por ejemplo
\codigo{REC-PCT-500-B}), y coordinar con Supervisión que las órdenes nuevas la usen.

\subseccion{Un material no existe en el sistema}
Pídeselo a Almacén o a Supervisión antes de empezar a capturar la receta. Tu rol no da
de alta materiales, y si sales del formulario a medias pierdes lo capturado.

\subseccion{¿Puedo crear una orden para probar mi receta?}
No. Pídesela a Supervisión o al Administrador.

\subseccion{¿La casilla \campo{Peligroso} restringe algo?}
No. Sólo cambia el aspecto de la pantalla del operario. El control real va en las
instrucciones y en el procedimiento de la planta.

\subseccion{¿Puedo definir el tipo de balanza o el instrumento?}
No. El sistema no conoce los instrumentos: \textbf{la báscula no está conectada} y el
operario teclea el peso. Lo único que controlas es el rango aceptado.

\subseccion{¿Cómo sé si una tolerancia es correcta?}
El sistema no opina. Contrasta el rango resultante contra la resolución real de la
balanza de piso y contra el expediente de desarrollo.
